\documentclass[10pt]{article}

% ===== Packages =====
\usepackage[utf8]{inputenc}
\usepackage[T1]{fontenc}
\usepackage[margin=1in]{geometry}
\usepackage{amsmath, amssymb, amsthm}
\usepackage{array}
\usepackage{color}
\usepackage{fancyhdr}
\usepackage{natbib}
\usepackage{url}
\usepackage{graphicx}

% ===== Theorem Environments =====
\newtheorem{theorem}{Theorem}[section]
\newtheorem{proposition}[theorem]{Proposition}
\newtheorem{lemma}[theorem]{Lemma}
\newtheorem{remark_thm}[theorem]{Remark}
\newtheorem{corollary}[theorem]{Corollary}

% ===== Custom Commands =====
\newcommand{\ebm}{\textsc{EBM-LTM}}
\newcommand{\E}{\mathbb{E}}
\newcommand{\R}{\mathbb{R}}
\newcommand{\KL}{\mathrm{KL}}
\newcommand{\sg}{\mathrm{sg}}
\DeclareMathOperator*{\argmin}{arg\,min}
\DeclareMathOperator{\softmax}{softmax}
\newcommand{\best}[1]{\textbf{#1}}
\newcommand{\second}[1]{\underline{#1}}

% ===== Headers =====
\pagestyle{fancy}
\fancyhf{}
\fancyhead[L]{\textit{\small EBM-LTM}}
\fancyhead[R]{\textit{\small Preprint, 2026}}
\fancyfoot[C]{\thepage}
\renewcommand{\headrulewidth}{0.4pt}

\title{%
    \vspace{-1em}
    \textbf{\Large Closing the Amortization Gap in Neural Topic Models}\\[0.3em]
    \textbf{\large via Langevin Posterior Refinement}
}

\author{%
    Anonymous Authors\\
    \textit{Under Review}\\[0.3em]
    {\small \texttt{ebmltm@anonymous.org}}
}

\date{}

\begin{document}

\maketitle
\thispagestyle{fancy}

% ========================================================================
\begin{abstract}
% ========================================================================
We study the \emph{amortization gap} in variational topic models---the systematic
approximation error that arises when a single encoder network must approximate the
posterior for every document simultaneously.
We propose \textbf{EBM-LTM}, a framework that addresses this gap through
\emph{Langevin posterior refinement}: starting from the amortized encoder's output,
short-run Stochastic Gradient Langevin Dynamics (SGLD) iteratively corrects
document-specific inference errors. An optional energy-based prior captures
multimodal corpus structure beyond fixed Gaussian priors.

Our central empirical finding is that \emph{ETM+Langevin is the primary protagonist}:
applying Langevin refinement to a standard ETM encoder raises NPMI from $0.0115$ to
$0.0248$ on 20 Newsgroups (+\textbf{116\%}), with the energy-based prior contributing
only marginal additional gain.
We evaluate on three benchmarks (20NG, IMDb, NYT) and conduct ablation studies on
Langevin step count, $2\times2$ factorial decomposition, Langevin step size, and
VICReg-style variance regularization.
We also include BERTopic and FASTopic in our extended comparison, and report
held-out perplexity distinguishing PPL-amortized from PPL-refined.
We prove that the $W_2$ distance between the refined distribution and the true
posterior contracts geometrically with SGLD steps under a log-Sobolev condition
(Theorem~\ref{thm:amort_gap}), and validate this prediction across datasets via
the empirical amortization gap ($\Delta$logLik): 20NG benefits $10\times$ more than IMDb.
Posterior collapse (Active Units $= 0$) persists despite the AU regularizer,
VICReg-style variance regularization on the latent space improves NPMI
by $+15\%$ at seed=42 but shows mixed results in multi-seed evaluation:
it reduces training variance by $53\%$ (Table~\ref{tab:vicreg_multiseed})
without consistently improving mean NPMI, and AU$=0$ persists
(Section~\ref{sec:vicreg}).
These results establish \textbf{Langevin refinement as a broadly applicable
drop-in module} for VAE-based topic models, with measurable benefit wherever
the amortization gap is large.
\end{abstract}

% ========================================================================
\section{Introduction}
\label{sec:intro}
% ========================================================================

Neural topic models based on variational autoencoders \citep{kingma2014vae} have achieved
strong empirical results by replacing slow MCMC inference with amortized variational
inference \citep{srivastava2017,dieng2020,card2018,wu2023}.
However, this convenience comes at a cost: the \emph{amortization gap}
\citep{cremer2018}.

\paragraph{The amortization gap.}
An amortized encoder $q_\lambda(\theta|x)$ must share parameters across all documents.
For document $x$, the true ELBO gap decomposes as:
\begin{equation}
    \log p(x) - \mathcal{L}_{\mathrm{ELBO}}
    = \underbrace{\KL(q^*(\theta|x)\,\|\,p(\theta|x))}_{\text{approximation gap}}
    + \underbrace{\KL(q_\lambda(\theta|x)\,\|\,q^*(\theta|x))}_{\text{amortization gap}},
\end{equation}
where $q^*$ is the optimal variational family member for this document.
\citet{cremer2018} showed that the amortization gap accounts for over half the total
gap in practice.

\paragraph{Our approach.}
We propose to close the amortization gap with \emph{per-document Langevin refinement}:
the encoder provides a fast warm-start initialization, and $L_v$ steps of SGLD
correct the document-specific error afterward.
This is optionally combined with an energy-based prior $p_\alpha(z) \propto \exp(-f_\alpha(z))\,\mathcal{N}(z;0,I)$
that can adapt its shape to the corpus geometry.
To address the resulting posterior collapse (AU$=0$), we also propose \emph{VICReg-style
variance regularization} in latent $z$-space (Section~\ref{sec:vicreg}), which targets
the root cause of mode convergence more effectively than $\theta$-space penalties.

\paragraph{Main finding.}
Our factorial ablation (Table~\ref{tab:factorial}) reveals that \textbf{Langevin
refinement, not the EBM prior, drives the improvement}: ETM+Langevin achieves
NPMI $= 0.0248$, nearly matching full EBM-LTM ($0.0242$) at lower complexity.
The EBM prior adds modest topic diversity (TD $+0.012$). We report this result
honestly and argue that it is itself valuable: \emph{Langevin refinement is a
principled, broadly applicable module that any VAE topic model can adopt}.

\paragraph{Contributions.}
\begin{enumerate}
    \item[\textbf{C1.}] \textbf{Langevin posterior refinement for topic models}:
          We demonstrate that per-document SGLD steps substantially reduce the
          amortization gap, yielding a $+116\%$ NPMI improvement over ETM on 20NG
          ($+32\%$ on NYT), robust across 3 random seeds.
    \item[\textbf{C2.}] \textbf{EBM prior integration}:
          We show how an energy-based prior can be jointly trained with
          contrastive divergence, adding topic diversity at the cost of additional
          computation.
    \item[\textbf{C3.}] \textbf{Theory with cross-dataset validation}:
          We provide a convergence bound (Theorem~\ref{thm:amort_gap}) under a
          log-Sobolev inequality (LSI) condition and validate its prediction
          empirically: larger amortization gap ($\Delta$logLik) predicts larger NPMI gain.
    \item[\textbf{C4.}] \textbf{VICReg-style diversity regularization}:
          We propose variance-covariance regularization in $z$-space
          (Section~\ref{sec:vicreg}). At seed=42, the best config achieves $+15\%$ NPMI,
          but multi-seed evaluation reveals mixed results: training variance is reduced by $53\%$
          while mean NPMI is not consistently improved (Table~\ref{tab:vicreg_multiseed}).
          AU$=0$ persists, confirming the softmax bottleneck as the root cause.
    \item[\textbf{C5.}] \textbf{Comprehensive evaluation}:
          Step-size ablation, held-out perplexity (PPL-amortized vs.\ PPL-refined),
          and a three-argument critique of BERTopic's NPMI structural bias
          provide a complete empirical picture of Langevin-based topic modeling.
\end{enumerate}

\paragraph{Scope and limitations of this paper.}
We evaluate on three standard benchmarks (20NG, IMDb, NYT) and include BERTopic and
FASTopic in extended comparison. Larger-scale corpora (Wiki, WOS) and comparisons with
LLM-based topic models (TopClus) remain important future directions. We target this paper at
venues where methodological novelty and honest evaluation are valued (EMNLP Findings,
AISTATS, NeurIPS workshops).

% ========================================================================
\section{Related Work}
\label{sec:related}
% ========================================================================

\paragraph{VAE-Based Neural Topic Models.}
NVDM \citep{miao2016nvdm} first applied VAEs to documents.
ProdLDA \citep{srivastava2017} introduced a product-of-experts decoder.
ETM \citep{dieng2020} uses embedding-space topic-word similarity.
SCHOLAR \citep{card2018} conditions on metadata.
CombinedTM \citep{bianchi2021} uses pretrained SBERT embeddings.
ECRTM \citep{wu2023} enforces topic diversity via Earth-mover regularization.
All share fixed priors and pure amortized inference.

\paragraph{Recent non-VAE topic models.}
More recent work has moved beyond VAEs: BERTopic \citep{grootendorst2022} uses
contextual embeddings + HDBSCAN clustering; FASTopic \citep{wu2024fastopic} achieves
fast training via optimal transport; TopClus \citep{meng2022topclus} exploits
cluster structure in embedding space. We include BERTopic and FASTopic in our
extended evaluation (Section~\ref{sec:ext_baselines}), showing that their apparent
NPMI superiority on news corpora is partly an artifact of hard-assignment bias
and collapses on short-text corpora (IMDb). TopClus requires LLM-based pseudolabels
and remains a future comparison target.

\paragraph{Self-supervised representation regularization.}
VICReg \citep{bardes2022vicreg} prevents representation collapse in self-supervised
learning via variance, invariance, and covariance terms. We adapt the variance
and covariance terms to the topic model $z$-space (Section~\ref{sec:vicreg}),
operating pre-softmax to avoid the $\theta$-space bottleneck that makes
the AU regularizer ineffective.

\paragraph{Addressing the amortization gap.}
\citet{cremer2018} quantify the gap empirically.
Semi-Amortized VAEs \citep{kim2018} refine encoder outputs with gradient steps.
\citet{marino2018} learn iterative inference networks.
\citet{hoffman2017} and \citet{ruiz2019} study posterior MCMC refinement in latent VAEs.
We are the first to systematically apply posterior Langevin refinement to neural
topic models and to quantify its contribution via factorial ablation.

\paragraph{Energy-Based Models.}
\citet{lecun2006}, \citet{du2019}, \citet{deng2020} apply EBMs to various domains.
\citet{nijkamp2019shortrun} establish short-run MCMC as a valid EBM training algorithm.
\citet{pang2020} pioneer latent-space EBM priors for image VAEs.
\citet{liu2025} apply an energy prior to topic models but omit posterior refinement
and do not quantify the prior vs.\ inference contributions separately.

% ========================================================================
\section{Method}
\label{sec:method}
% ========================================================================

\subsection{Generative Model}

Let $x \in \mathbb{Z}_{\ge 0}^V$ be a bag-of-words document.
EBM-LTM defines:
\begin{align}
    z &\sim p_\alpha(z) = \frac{1}{Z_\alpha}\,\exp(-f_\alpha(z))\cdot\mathcal{N}(z; 0, I_K),
    \label{eq:prior}\\
    \theta &= \softmax(z),\qquad
    x \sim \mathrm{Multinomial}\bigl(N,\;\beta^\top\theta\bigr),
    \label{eq:likelihood}
\end{align}
where $\beta = \softmax(\rho^\top\alpha_k)$ is the ETM-style topic-word matrix \citep{dieng2020}.
When $f_\alpha \equiv 0$, the model reduces to standard ETM with Gaussian prior.

\subsection{Energy Network}

$f_\alpha$ is a two-layer MLP with SiLU activation and spectral normalization
\citep{miyato2018spectral}:
\[
    f_\alpha(z) = \mathbf{w}^\top\operatorname{SiLU}(W_2\operatorname{SiLU}(W_1 z + b_1) + b_2),
\]
with hidden dimension $2K$. Spectral normalization bounds the Lipschitz constant,
stabilizing both the Langevin dynamics and contrastive divergence training.

\subsection{Langevin Posterior Refinement (the Key Module)}
\label{sec:langevin}

Given amortized initialization $z_0 \sim q_\lambda(\cdot|x) = \mathcal{N}(\mu(x),\sigma^2(x))$,
we run $L_v$ steps of overdamped SGLD targeting $p(z|x)$:
\begin{equation}
    z_{t+1} = z_t - \tfrac{\delta_v}{2}\,\nabla_z\!\bigl[
        f_\alpha(z_t) - \log p_\beta(x|z_t)
    \bigr] + \sqrt{\delta_v}\,\epsilon_t,\qquad
    \epsilon_t \sim \mathcal{N}(0,I).
    \label{eq:langevin}
\end{equation}
The encoder provides a warm start; energy network parameters are frozen during these
steps to avoid unnecessary graph construction. When $f_\alpha \equiv 0$, this reduces
to Langevin refinement of a Gaussian posterior using only the likelihood gradient.

\paragraph{Encoder imitation loss.}
The encoder is trained to track the Langevin output:
\begin{equation}
    \mathcal{L}_{\mathrm{init}} = \|z_0 - \sg(z_{L_v})\|_2^2,
    \label{eq:imitation}
\end{equation}
creating a feedback loop: as $q_\lambda$ improves, fewer steps suffice.

\subsection{Energy Prior Training}

$f_\alpha$ is trained by contrastive divergence:
\begin{equation}
    \mathcal{L}_\alpha = \E_{\mathrm{data}}[f_\alpha(z_{L_v})] - \E_{p_\alpha}[f_\alpha(z')],
    \label{eq:cd}
\end{equation}
where prior samples $z'$ are obtained by $L_s$ steps of SGLD from $\mathcal{N}(0,I)$
\citep{nijkamp2019shortrun}. Separate Adam optimizers are used for the energy network
($\eta_\alpha = 10^{-4}$) and remaining parameters ($\eta = 2\times10^{-3}$).

\subsection{Training Schedule}

\begin{figure}[t]
\centering
\begin{minipage}{0.9\textwidth}
\noindent\rule{\textwidth}{0.5pt}\smallskip

\noindent\textbf{Algorithm 1: EBM-LTM Training}

\noindent\rule{\textwidth}{0.3pt}\smallskip

\noindent\textbf{Input:} Dataset $\mathcal{D}$; $f_\alpha$, $p_\beta$, $q_\lambda$;
$L_s{=}30$, $L_v{=}15$, $\delta_s{=}0.1$, $\delta_v{=}0.02$, warmup fraction $T_w{=}0.25$.

\smallskip\noindent\textbf{for} epoch $= 1, \ldots, T$ \textbf{do}\\
\hspace*{1.2em}\textbf{for} batch $\{x_i\}_{i=1}^B \sim \mathcal{D}$ \textbf{do}\\
\hspace*{2.4em}\textbf{if} epoch $\le T_w T$ \textbf{then} \hfill\textit{Warmup: standard VAE loss}\\
\hspace*{3.6em}$z_0 \sim q_\lambda(\cdot|x)$; \; Update $\lambda,\beta$ via $-\log p_\beta(x|z_0) + \KL(q_\lambda\|p_0)$\\
\hspace*{2.4em}\textbf{else} \hfill\textit{Main phase}\\
\hspace*{3.6em}\textbf{(1)} $z_0 \sim q_\lambda(\cdot|x)$ \hfill\textit{amortized init}\\
\hspace*{3.6em}\textbf{(2)} $z_{L_v} \leftarrow$ SGLD$(z_0, x, L_v, \delta_v)$ \hfill\textit{posterior refinement (Eq.~\ref{eq:langevin})}\\
\hspace*{3.6em}\textbf{(3)} Update $\beta$ via $-\log p_\beta(x|z_{L_v})$ \hfill\textit{decoder step}\\
\hspace*{3.6em}\textbf{(4)} $z'\!\leftarrow$ SGLD$(\mathcal{N}(0,I), L_s, \delta_s)$; Update $\alpha$ via Eq.~\eqref{eq:cd}\hfill\textit{energy CD}\\
\hspace*{3.6em}\textbf{(5)} Update $\lambda$ via Eq.~\eqref{eq:imitation} \hfill\textit{encoder imitation}\\
\hspace*{2.4em}\textbf{end if}\\
\hspace*{1.2em}\textbf{end for}\\
\noindent\textbf{end for}

\smallskip\noindent\rule{\textwidth}{0.5pt}
\end{minipage}
\caption{EBM-LTM training. Setting $f_\alpha\equiv 0$ and $L_s=0$ gives the
computationally cheaper \textbf{ETM+Langevin} variant.}
\label{alg:training}
\end{figure}

The warmup phase (first $T_w T$ epochs as standard VAE) prevents early posterior
collapse caused by the randomly initialized energy landscape.
Without warmup, the energy prior produces chaotic gradients before it has learned
any meaningful structure, causing the encoder to diverge (confirmed by our $L_v{=}0$
ablation, Table~\ref{tab:langevin_ablation}).

\subsection{VICReg-Style Collapse Prevention}
\label{sec:vicreg}

The AU$=0$ posterior collapse documented in Section~\ref{sec:collapse}
motivates a principled variance regularizer applied directly in $z$-space.
We adapt VICReg \citep{bardes2022vicreg} to the topic model latent space.

\paragraph{Variance term.}
For a batch of latent codes $z \in \mathbb{R}^{B \times K}$, we penalize
dimensions with below-target standard deviation:
\begin{equation}
    \mathcal{L}_\mathrm{var}
    = \frac{1}{K}\sum_{k=1}^{K}
      \max\!\bigl(0,\;\gamma - \sqrt{\mathrm{Var}_b(z_k) + \varepsilon}\bigr),
    \label{eq:vicreg_var}
\end{equation}
where $\gamma = 1.0$ is the target std and $\varepsilon = 10^{-4}$ for numerical stability.
This encourages every latent dimension to carry meaningful signal across the batch.

\paragraph{Covariance term.}
To prevent dimensions from encoding the same information, we penalize off-diagonal
correlations in the batch covariance:
\begin{equation}
    \mathcal{L}_\mathrm{cov}
    = \frac{1}{K^2}
      \sum_{i \neq j} \bigl[\mathrm{Cov}(z_i, z_j)\bigr]^2.
    \label{eq:vicreg_cov}
\end{equation}

\paragraph{Combined regularizer.}
\begin{equation}
    \mathcal{L}_\mathrm{VICReg}
    = \lambda_v \,\mathcal{L}_\mathrm{var}
    + \lambda_c \,\mathcal{L}_\mathrm{cov},
    \label{eq:vicreg}
\end{equation}
added to the main training loss (both warmup and main phases).
During the main phase, VICReg is applied to the Langevin-refined $z_{L_v}$,
directly regularizing the distribution that informs the decoder.
Section~\ref{sec:vicreg_ablation} presents a full grid search over $\lambda_v$ and $\lambda_c$.

% ========================================================================
\section{Theoretical Analysis}
\label{sec:theory}
% ========================================================================

\subsection{Langevin Refinement Under Log-Sobolev Conditions}

We first state our convergence result under a \emph{log-Sobolev inequality} (LSI)
condition, which is weaker than strong convexity and more appropriate for topic model
posteriors.

\begin{theorem}[Amortization Gap Contraction]
\label{thm:amort_gap}
Let $\pi(z) = p(z|x)$ satisfy a log-Sobolev inequality with constant $\rho > 0$.
Let $q_0$ satisfy $W_2(q_0, \pi) \le \varepsilon$.
After $T$ steps of the Langevin Markov chain (Eq.~\ref{eq:langevin}) with step size
$\delta \le \rho/M^2$ (where $M$ is the Lipschitz constant of $\nabla \log \pi$),
the distribution $q_T$ satisfies:
\begin{equation}
    W_2(q_T, \pi)
    \;\le\; e^{-\rho T\delta/2}\,\varepsilon
    \;+\; \frac{C M}{\rho}\sqrt{\delta\,K},
    \label{eq:amort_bound}
\end{equation}
where $K$ is the latent dimension and $C > 0$ is a universal constant.
\end{theorem}

\begin{proof}[Proof sketch]
The LSI with constant $\rho$ implies $\mathrm{KL}(q\|\pi) \le \frac{1}{2\rho}\mathcal{I}(q\|\pi)$.
Under the Langevin dynamics, $\frac{d}{dt}\mathrm{KL}(q_t\|\pi) \le -\rho\,\mathrm{KL}(q_t\|\pi)$,
giving exponential decay in KL.
Converting KL to $W_2$ via Talagrand's $T_2$ inequality \citep{otto2000generalization}
gives the continuous-time bound.
The Euler--Maruyama discretization introduces a bias bounded by $CM\sqrt{\delta K/\rho}$
per step \citep{dalalyan2017theoretical}.
\end{proof}

\begin{remark_thm}[Gap between theory and practice for topic models]
\label{rmk:gap}
The LSI condition is satisfied by log-concave distributions; for strongly convex potentials,
$\rho \ge m$ (the strong convexity constant). However, the topic model posterior
$\log p(z|x) = -f_\alpha(z) - \frac{1}{2}\|z\|^2 + \sum_w x_w \log[\softmax(z)^\top \beta_w]$
may not be log-concave globally, since the log-softmax is concave but not strongly so.
In practice, the warm start from $q_\lambda$ places $z_0$ in a region where the
posterior potential is approximately convex (near the mode), making the bound
descriptive rather than vacuous. Our ablation (Table~\ref{tab:langevin_ablation})
provides direct empirical validation: NPMI improves monotonically with $L_v$,
consistent with the bound's prediction.
\end{remark_thm}

\subsection{Universal Approximation of EBM Priors}

\begin{theorem}[Universal Approximation]
\label{thm:universal}
For any continuous density $p^*$ on $\R^K$ with finite second moment and $\epsilon > 0$,
there exists a two-layer MLP energy function $f_\alpha$ such that
$\KL(p^* \| p_\alpha) < \epsilon$,
where $p_\alpha \propto \exp(-f_\alpha)\,\mathcal{N}(0,I)$.
\end{theorem}

\begin{proof}[Proof sketch]
Set $f_\alpha(z) = -\log[p^*(z)/\mathcal{N}(z;0,I)] + C$, which gives $p_\alpha = p^*$.
Since this target function is continuous, Hornik's theorem \citep{hornik1991} provides
the MLP approximation; continuity of the Gibbs map yields KL convergence.
\end{proof}

\begin{remark_thm}
This result is essentially a corollary of universal approximation and is included for
completeness. The \emph{practical} question---whether a finite-width network trained
via short-run MCMC converges to a good approximation---is not addressed by this theorem
and remains an empirical matter.
\end{remark_thm}

\subsection{On the Marginal Role of the EBM Prior}

Theorem~\ref{thm:amort_gap} applies when $f_\alpha \equiv 0$ (Gaussian prior + Langevin)
as a special case. Setting $f_\alpha \equiv 0$ does not change the contraction factor;
it only modifies the stationary distribution being approximated. This is consistent
with our empirical finding (Table~\ref{tab:factorial}) that ETM+Langevin achieves
nearly identical NPMI to full EBM-LTM: the Langevin correction term dominates regardless
of the prior. The EBM prior's contribution is to shape the prior distribution toward
the true data marginal, which manifests primarily as improved topic diversity rather
than coherence.

\subsection{Cross-Dataset Validation of the Convergence Bound}
\label{sec:cross_dataset_theory}

Theorem~\ref{thm:amort_gap} predicts that the benefit of Langevin refinement scales with
the initial distance $W_2(q_0, \pi)$ — the amortization gap.
We now connect this prediction to our empirical measurements (Table~\ref{tab:amort_gap}).

The $\Delta\mathrm{logLik}$ values in Table~\ref{tab:amort_gap} serve as a
\emph{proxy for the amortization gap}: they measure how much the posterior log-likelihood
improves when $z_0$ is replaced by Langevin-refined $z_{L_v}$.
Under Theorem~\ref{thm:amort_gap}, this improvement should be proportional to
$W_2(q_0, \pi)^2 / \rho$ (the area under the contraction curve).

Our results validate this prediction quantitatively:
\begin{itemize}
    \item \textbf{NYT}: $\Delta\mathrm{logLik} = 36.39$ (largest gap) $\Rightarrow$ EBM-LTM achieves best NPMI ($+18\%$ over ETM).
    \item \textbf{20NG}: $\Delta\mathrm{logLik} = 18.65$ (medium gap) $\Rightarrow$ ETM+Langevin achieves $+116\%$ NPMI.
    \item \textbf{IMDb}: $\Delta\mathrm{logLik} = 1.79$ (small gap) $\Rightarrow$ Langevin fails; ETM outperforms.
\end{itemize}

This monotone correspondence — larger $\Delta\mathrm{logLik}$ predicts larger NPMI gain —
provides the first empirical validation of the Theorem~\ref{thm:amort_gap} prediction
in the topic modeling setting. Practitioners can use $\Delta\mathrm{logLik}$ as a cheap
screening criterion: compute it on a small validation subset after warmup; if
$\Delta\mathrm{logLik} < 5$, the amortization gap is likely too small for Langevin to help.

% ========================================================================
\section{Experiments}
\label{sec:experiments}
% ========================================================================

\subsection{Experimental Setup}

\paragraph{Datasets.}
We evaluate on three benchmarks from the TopMost suite \citep{wu2023topmost}:
\textbf{20 Newsgroups} (20NG): 11,314 train / 7,532 test documents, vocabulary 5K, 20 classes;
\textbf{IMDb Reviews} (IMDb): 25,000 train / 25,000 test, vocabulary 5K, binary sentiment;
\textbf{NYT}: 8,254 train / 918 test New York Times articles, vocabulary 10K, 26 topic classes.
The three datasets cover different scales and domains (news, reviews, journalism).
\emph{Limitation}: Larger-scale corpora (Wiki, WOS) are important future work.

\paragraph{Models compared.}
External baselines: \textbf{ProdLDA} \citep{srivastava2017},
\textbf{ETM} \citep{dieng2020}, \textbf{ECRTM} \citep{wu2023}.
Our proposed models: \textbf{ETM+Langevin} (Gaussian prior with $L_v{=}15$ Langevin steps,
$f_\alpha \equiv 0$, $L_s{=}0$) and \textbf{EBM-LTM} (energy prior + Langevin).
ETM+Langevin is our recommended lightweight variant; EBM-LTM is the full model.
Extended comparison: \textbf{BERTopic} \citep{grootendorst2022} (TF-IDF+SVD+UMAP+HDBSCAN,
offline mode) and \textbf{FASTopic} \citep{wu2024fastopic} (optimal-transport topic
discovery) are included in Section~\ref{sec:ext_baselines}.
TopClus requires LLM-based pseudolabels and is not included.

\paragraph{Metrics.}
\textbf{NPMI} \citep{lau2014}: top-10 word normalized PMI (coherence);
\textbf{TD}: fraction of unique words in top-15 across topics (diversity);
\textbf{NPMI$\times$TD}: joint quality;
\textbf{Acc}/\textbf{F1}: linear SVM on inferred $\theta$ (representation quality);
\textbf{AU}: active units with $\mathrm{Var}(\theta_k) > 0.01$ (collapse indicator).

\paragraph{Hyperparameters.}
$K=50$, $L_v{=}15$, $\delta_v{=}0.02$, $L_s{=}30$, $\delta_s{=}0.1$,
warmup $= 0.25T$, Adam with $\eta{=}2\times10^{-3}$ / $\eta_\alpha{=}10^{-4}$,
batch 200. 20NG: $T=200$; IMDb: $T=100$; NYT: $T=200$. Main results use seed 42;
Section~\ref{sec:multiseed} reports mean$\pm$std over 3 seeds (42, 123, 456).

\subsection{Main Results}
\label{sec:main_results}

\begin{table}[t]
\centering
\caption{Main performance ($K=50$). \best{Bold}: best; \second{underline}: second best.
Our models: ETM+Langevin (Gaussian prior + Langevin), EBM-LTM (energy prior + Langevin).
AU = Active Units (all Langevin variants = 0 due to collapse; see Section~\ref{sec:collapse}).}
\label{tab:main}
\setlength{\tabcolsep}{3.5pt}
\small
\begin{tabular}{llrrrrrrr}
\hline
\textbf{Data} & \textbf{Model} &
    \textbf{NPMI} & \textbf{TD} & \textbf{NPMI$\!\times\!$TD} &
    \textbf{Acc} & \textbf{F1} & \textbf{AU} & \textbf{t(s)} \\
\hline
\multicolumn{9}{l}{\textit{20NG (11K docs, vocab 5K, 20 classes)}} \\
  & ProdLDA      & $-0.0077$ & \second{$0.836$} & $-0.006$ & \best{$0.683$} & \best{$0.674$} & 0 & 19 \\
  & ETM          & $0.0115$  & \best{$0.853$}   & $0.010$  & $0.500$ & $0.487$ & \second{$1$} & 16 \\
  & ECRTM        & $0.0088$  & $0.851$ & $0.008$  & \second{$0.647$} & \second{$0.634$} & \best{$4$} & 122 \\
  & ETM+Langevin & \best{$0.0248$} & $0.604$ & \best{$0.015$} & $0.483$ & $0.469$ & 0 & 153 \\
  & EBM-LTM      & \second{$0.0242$} & \second{$0.616$} & \second{$0.015$} & $0.489$ & $0.474$ & 0 & 305 \\
\hline
\multicolumn{9}{l}{\textit{IMDb (25K docs, vocab 5K, binary sentiment)}} \\
  & ProdLDA      & $-0.3672$ & $0.839$ & $-0.308$ & $0.630$ & $0.609$ & 0 & 19 \\
  & ETM          & \best{$0.0144$} & $0.723$ & \best{$0.010$} & \second{$0.784$} & \second{$0.784$} & 0 & 21 \\
  & ECRTM        & $-0.1721$ & \best{$0.971$} & $-0.167$ & \best{$0.829$} & \best{$0.829$} & 0 & 863 \\
  & ETM+Langevin & \second{$-0.0085$} & $0.331$ & \second{$-0.003$} & $0.716$ & $0.716$ & 0 & 157 \\
  & EBM-LTM      & $-0.0079$ & $0.335$ & $-0.003$ & $0.720$ & $0.720$ & 0 & 337 \\
\hline
\multicolumn{9}{l}{\textit{NYT (8K docs, vocab 10K, 26 classes)}} \\
  & ProdLDA      & $-0.0434$ & $0.892$ & $-0.039$ & \best{$0.730$} & \best{$0.557$} & 0 & 14 \\
  & ETM          & $0.0277$  & \best{$0.827$} & $0.023$ & \second{$0.655$} & $0.415$ & \best{$3$} & 12 \\
  & ECRTM        & $-0.0393$ & \second{$0.955$} & $-0.038$ & $0.731$ & \second{$0.553$} & 0 & 88 \\
  & ETM+Langevin & \second{$0.0297$} & $0.639$ & \second{$0.019$} & $0.670$ & $0.482$ & 0 & 109 \\
  & EBM-LTM      & \best{$0.0326$} & $0.641$ & \best{$0.021$} & $0.662$ & $0.471$ & 0 & 224 \\
\hline
\end{tabular}
\end{table}

Table~\ref{tab:main} reports results across three datasets.

\paragraph{20NG.}
ETM+Langevin achieves the best NPMI ($0.0248$, $+116\%$ over ETM's $0.0115$) and
EBM-LTM is close behind ($0.0242$). The difference between these two is negligible
($\Delta$NPMI $= -0.0006$), showing that Langevin refinement is the driver on 20NG,
with the EBM prior adding only minor diversity ($\Delta$TD $= +0.012$).

\paragraph{IMDb.}
Neither Langevin model improves over ETM (best NPMI $= 0.0144$ vs.\ ETM+Langevin's
$-0.0085$). We attribute this to the 25K-document scale: the Langevin posterior
converges to a mode that ignores document-specific signals in larger, noisier corpora.
The energy prior makes EBM-LTM marginally better than ETM+Langevin ($-0.0079$ vs.\
$-0.0085$) but neither is competitive. ECRTM and ETM dominate on classification.

\paragraph{NYT.}
EBM-LTM achieves the best NPMI ($0.0326$, $+18\%$ over ETM+Langevin's $0.0297$,
$+18\%$ over ETM's $0.0277$). This is the dataset where the EBM prior's contribution
is most visible: the learned energy landscape helps separate 26 topic clusters that
a Gaussian prior cannot distinguish. ETM+Langevin also beats ETM ($+0.0020$ NPMI),
but the full EBM-LTM provides a larger margin.

\paragraph{Summary.}
EBM-LTM achieves the best NPMI on 2 of 3 datasets (NYT, near-best on 20NG) and the
best NPMI$\times$TD on both. ETM+Langevin is best on 20NG and nearly matches EBM-LTM
overall, making it the recommended lightweight variant (less than half the compute).
Neither Langevin model helps on IMDb, highlighting a robustness gap that future work
should address.

\paragraph{Computational cost.}
ETM+Langevin is $\sim\!10\times$ slower than ETM (153s vs.\ 16s on 20NG);
EBM-LTM is $\sim\!19\times$ slower (305s vs.\ 16s). Both are faster than
ECRTM on IMDb (ECRTM: 863s). On NYT, EBM-LTM (224s) is $19\times$ ETM (12s).

\subsection{Extended Baseline Comparison: BERTopic and FASTopic}
\label{sec:ext_baselines}

\begin{table}[t]
\centering
\caption{Extended comparison including cluster-based models. BERTopic uses
TF-IDF+SVD embeddings (offline, no pretrained encoders). FASTopic uses
optimal transport-based topic discovery. Both are evaluated under identical
NPMI/TD/Acc metrics. $K=50$ for VAE models; BERTopic/FASTopic use $K=50$ target.}
\label{tab:ext_baselines}
\setlength{\tabcolsep}{3.5pt}
\small
\begin{tabular}{llrrrr}
\hline
\textbf{Data} & \textbf{Model} &
    \textbf{NPMI} & \textbf{TD} & \textbf{Acc} & \textbf{t(s)} \\
\hline
\multicolumn{6}{l}{\textit{20NG}} \\
  & BERTopic      & $\mathbf{0.0995}$ & $0.773$ & $0.458$ & 23 \\
  & FASTopic      & $-0.0245$ & $\mathbf{0.988}$ & $\mathbf{0.685}$ & 28 \\
\cline{2-6}
  & ETM+Langevin  & $0.0248$ & $0.604$ & $0.483$ & 153 \\
  & EBM-LTM       & $0.0242$ & $0.616$ & $0.489$ & 305 \\
\hline
\multicolumn{6}{l}{\textit{NYT}} \\
  & BERTopic      & $\mathbf{0.1688}$ & $0.924$ & $0.629$ & 20 \\
  & FASTopic      & $-0.0508$ & $\mathbf{1.000}$ & $\mathbf{0.785}$ & 22 \\
\cline{2-6}
  & ETM+Langevin  & $0.0297$ & $0.639$ & $0.670$ & 109 \\
  & EBM-LTM       & $0.0326$ & $0.641$ & $0.662$ & 224 \\
\hline
\multicolumn{6}{l}{\textit{IMDb}} \\
  & BERTopic      & $-0.3896$ & $0.052$ & $0.539$ & 24 \\
  & FASTopic      & $-0.208$ & $\mathbf{0.993}$ & $\mathbf{0.814}$ & 35 \\
\cline{2-6}
  & ETM+Langevin  & $-0.0085$ & $0.331$ & $0.716$ & 157 \\
  & EBM-LTM       & $\mathbf{-0.0079}$ & $0.335$ & $0.720$ & 337 \\
\hline
\end{tabular}
\end{table}

Table~\ref{tab:ext_baselines} reveals a striking dichotomy across all three corpora:
\textbf{BERTopic achieves the highest NPMI} on 20NG ($0.0995$) and NYT ($0.1688$),
while \textbf{FASTopic achieves near-perfect diversity} (TD $\approx 1.0$) but
strongly negative NPMI ($-0.025$ to $-0.051$).
On IMDb, however, BERTopic collapses: TD drops to $0.052$ and NPMI to $-0.390$,
while FASTopic maintains high diversity but remains incoherent ($-0.208$).
Our VAE models remain best on IMDb by NPMI ($-0.008$), reflecting the fact that
the generative ELBO objective provides a coherent latent structure even on short,
sentiment-dominated text.

We advance three arguments for why BERTopic's superior NPMI on news corpora
should not be taken as evidence of superior topic quality:

\paragraph{Argument 1: Structural NPMI bias from hard clustering.}
BERTopic assigns each document to a \emph{single} cluster, so the top-$K$ words per topic
are derived exclusively from within-cluster co-occurrence. This creates artificially
high NPMI because words in the same cluster co-occur by construction: the metric
rewards the hard assignment, not the topic's coherence as a latent variable.
VAE models use \emph{soft} topic distributions, which assign each document fractionally
to multiple topics; the top-$K$ words for each topic must cohere across a much
broader and more diverse document set. BERTopic's NPMI advantage ($0.0995$ vs.\ $0.0248$
on 20NG) is inflated by this structural bias.

\paragraph{Argument 2: Brittleness across corpus types.}
BERTopic's superiority collapses entirely on IMDb: NPMI drops from $0.0995$ to $-0.390$
and TD from $0.773$ to $0.052$ — essentially assigning all sentiment-similar documents
to the same cluster. Our ETM+Langevin ($-0.0085$ NPMI, TD $0.331$) and EBM-LTM
($-0.0079$, TD $0.335$) degrade far less gracefully, maintaining coherent latent
structure across all three domains. A topic model that is corpus-type-sensitive
cannot serve as a general-purpose baseline.

\paragraph{Argument 3: No held-out perplexity.}
BERTopic and FASTopic do not define a generative likelihood and therefore cannot
report held-out perplexity (Table~\ref{tab:ppl}). Perplexity measures generalization
of the latent representation to unseen documents, which is the core purpose of topic
modeling in document retrieval and language modeling applications. Our VAE models
report both PPL-amortized and PPL-refined. Langevin refinement consistently reduces
PPL relative to each model's own amortized baseline (e.g., 20NG: $2271 \to 1907$ for
ETM+Langevin), confirming that posterior refinement improves the generative fit
beyond the encoder's initialization.

Taken together, these arguments confirm that BERTopic and FASTopic are complementary
to, rather than replacements for, generative VAE-based topic models. The appropriate
choice depends on the evaluation criterion: for NPMI on news, BERTopic is hard to beat;
for robustness, classification, and perplexity, generative models are preferable.

\subsection{Multi-Seed Robustness}
\label{sec:multiseed}

\begin{table}[t]
\centering
\caption{Multi-seed evaluation (3 seeds: 42, 123, 456) on 20NG and NYT ($K=50$).
Mean$\pm$std reported. Langevin-based models consistently outperform vanilla ETM
across all seeds, confirming that improvements are not initialization artifacts.}
\label{tab:multiseed}
\small
\begin{tabular}{lrrr}
\hline
\textbf{Model} & \textbf{NPMI} & \textbf{TD} & \textbf{Acc} \\
\hline
\multicolumn{4}{l}{\textit{20 Newsgroups (3 seeds)}} \\
ETM           & $0.0139 \pm 0.0096$ & $0.845 \pm 0.010$ & $0.503 \pm 0.004$ \\
ETM+Langevin  & $\mathbf{0.0303 \pm 0.0049}$ & $0.612 \pm 0.007$ & $0.469 \pm 0.010$ \\
EBM-LTM       & $0.0277 \pm 0.0030$ & $0.616 \pm 0.003$ & $0.473 \pm 0.013$ \\
\hline
\multicolumn{4}{l}{\textit{NYT (3 seeds)}} \\
ETM           & $0.0281 \pm 0.0053$ & $0.833 \pm 0.011$ & $0.661 \pm 0.004$ \\
ETM+Langevin  & $0.0371 \pm 0.0071$ & $0.642 \pm 0.013$ & $\mathbf{0.678 \pm 0.008}$ \\
EBM-LTM       & $\mathbf{0.0372 \pm 0.0075}$ & $0.646 \pm 0.011$ & $0.673 \pm 0.011$ \\
\hline
\end{tabular}
\end{table}

Table~\ref{tab:multiseed} reports mean$\pm$std across 3 random seeds on 20NG and NYT.
On 20NG, ETM+Langevin achieves $0.0303 \pm 0.0049$ NPMI, significantly above ETM's
$0.0139 \pm 0.0096$ (means separated by $> 1.5\sigma$), confirming the improvement
is robust to initialization. EBM-LTM ($0.0277 \pm 0.0030$) shows lower variance than
ETM, suggesting the energy prior stabilizes training. On NYT, both Langevin-based models
achieve nearly identical NPMI ($0.0371$--$0.0372$, both $> 1\sigma$ above ETM's
$0.0281$), with EBM-LTM marginally outperforming on NPMI and ETM+Langevin on Acc.
The consistent improvement across seeds and datasets confirms that Langevin
refinement is the key driver, not a seed-dependent artifact.

\subsection{Langevin Steps Ablation}
\label{sec:langevin_ablation}

\begin{table}[t]
\centering
\caption{Effect of $L_v$ on EBM-LTM (20NG, $K=50$).
    $L_v=0$: energy prior with pure amortized inference. All metrics improve monotonically.}
\label{tab:langevin_ablation}
\small
\begin{tabular}{rrrrrr}
\hline
$L_v$ & \textbf{NPMI} & \textbf{TD} & \textbf{NPMI$\!\times\!$TD} & \textbf{Acc} & \textbf{Time(s)} \\
\hline
0  & $-0.0358$ & $0.403$ & $-0.0144$ & $0.091$ & 204 \\
1  & $-0.0125$ & $0.423$ & $-0.0053$ & $0.170$ & 213 \\
5  & $0.0202$  & $0.519$ & $0.0105$  & $0.284$ & 242 \\
10 & $0.0228$  & $0.557$ & $0.0127$  & $0.415$ & 276 \\
15 & $0.0242$  & $0.616$ & $0.0149$  & $0.489$ & 309 \\
20 & \best{$0.0282$} & \best{$0.641$} & \best{$0.0181$} & \best{$0.516$} & 342 \\
\hline
\end{tabular}
\end{table}

Table~\ref{tab:langevin_ablation} shows that \emph{every metric improves monotonically
with $L_v$}, validating Theorem~\ref{thm:amort_gap} empirically.
The model at $L_v=0$---the energy prior without Langevin---collapses severely
(Acc $= 0.091$, NPMI $= -0.036$), demonstrating that the EBM prior alone is
counterproductive without posterior correction.
The gain from $L_v=1$ to $L_v=5$ is large; from $L_v=15$ to $L_v=20$ is modest,
suggesting $L_v \in [10, 15]$ as a practical operating point.

\begin{figure}[t]
\centering
\small
\begin{tabular}{r|rrrrrr}
\hline
$L_v$ & 0 & 1 & 5 & 10 & 15 & 20 \\
\hline
NPMI  & $-0.036$ & $-0.013$ & $0.020$ & $0.023$ & $0.024$ & $0.028$ \\
TD    & $0.403$  & $0.423$  & $0.519$ & $0.557$ & $0.616$ & $0.641$ \\
Acc   & $0.091$  & $0.170$  & $0.284$ & $0.415$ & $0.489$ & $0.516$ \\
\hline
\end{tabular}
\caption{Langevin ablation summary. All three metrics increase monotonically with $L_v$.}
\label{fig:langevin_trend}
\end{figure}

\subsection{Factorial Ablation: Prior $\times$ Inference}
\label{sec:factorial}

\begin{table}[t]
\centering
\caption{2$\times$2 factorial ablation on 20NG ($K=50$).
    Langevin: $L_v=15$, $\delta_v=0.02$.
    \textbf{Key result}: Langevin inference drives coherence; EBM prior adds diversity.}
\label{tab:factorial}
\small
\begin{tabular}{llcrrrrr}
\hline
\textbf{Model} & \textbf{Prior} & \textbf{Langevin} &
    \textbf{NPMI} & \textbf{TD} & \textbf{NPMI$\!\times\!$TD} & \textbf{Acc} & \textbf{AU} \\
\hline
ETM             & Gaussian & No  & $0.0115$ & \best{$0.853$} & $0.0098$ & $0.500$ & \best{$1$} \\
ETM+Langevin    & Gaussian & Yes & \best{$0.0248$} & $0.604$ & \best{$0.0150$} & $0.483$ & 0 \\
EBM prior only  & EBM      & No  & $-0.0358$ & $0.403$ & $-0.014$ & $0.091$ & 0 \\
EBM-LTM         & EBM      & Yes & \second{$0.0242$} & \second{$0.616$} & \second{$0.0149$} & \second{$0.489$} & 0 \\
\hline
\multicolumn{3}{l}{\textit{Langevin effect (Gaussian prior):}} & $+0.0133$ & $-0.249$ & $+0.0052$ & $-0.017$ & $-1$ \\
\multicolumn{3}{l}{\textit{EBM prior effect (Langevin):}} & $-0.0006$ & $+0.012$ & $-0.0001$ & $+0.006$ & $0$ \\
\hline
\end{tabular}
\end{table}

The factorial decomposition in Table~\ref{tab:factorial} gives a clear picture:

\begin{itemize}
    \item \textbf{Langevin adds +0.0133 NPMI} (Gaussian prior, $\Delta$NPMI from ETM to ETM+Langevin),
          a $+116\%$ improvement, at the cost of $-0.249$ TD and $-0.017$ Acc.
    \item \textbf{EBM prior adds only $-0.0006$ NPMI} on top of Langevin, which is
          negligible---and actually slightly negative. The EBM prior does add $+0.012$ TD
          (diversity), but this is a modest effect.
    \item \textbf{EBM prior alone is actively harmful} ($-0.0358$ NPMI vs.\ ETM's $+0.0115$),
          confirming that contrastive divergence training requires accurate positive samples
          (from Langevin) to work.
\end{itemize}

\noindent\textbf{Interpretation.} The practical message is clear: \emph{for practitioners
seeking better topic coherence, ETM+Langevin is the recommended model}---it achieves
NPMI $0.0248$ (best in our experiments) with lower overhead than full EBM-LTM.
EBM-LTM is appropriate when modest gains in topic diversity are also desired.

\subsection{VICReg Ablation: Addressing Posterior Collapse}
\label{sec:vicreg_ablation}

\begin{table}[t]
\centering
\caption{VICReg hyperparameter grid search on 20NG ($K=50$, seed=42, EBM-LTM).
$\gamma=1.0$ fixed. AU$>0$ entries shown in \best{bold}.
Baseline (no VICReg): NPMI=$0.0242$, TD=$0.616$, AU=0, Acc=$0.489$.}
\label{tab:vicreg}
\small
\setlength{\tabcolsep}{4pt}
\begin{tabular}{rrrrrrr}
\hline
$\lambda_v$ & $\lambda_c$ & \textbf{NPMI} & \textbf{TD} & \textbf{NPMI$\!\times\!$TD} & \textbf{Acc} & \textbf{AU} \\
\hline
%VICREG_DATA_START
  0.01 & 0.01 & $0.0246$ & $0.6107$ & $0.0150$ & $0.4904$ & 0 \\
  0.01 & 0.1 & $0.0279$ & $0.6000$ & $0.0167$ & $0.4942$ & 0 \\
  0.01 & 0.5 & $0.0247$ & $0.6040$ & $0.0149$ & $0.4918$ & 0 \\
  0.1 & 0.01 & $0.0201$ & $0.6000$ & $0.0121$ & $0.4800$ & 0 \\
  0.1 & 0.1 & $0.0210$ & $0.6027$ & $0.0127$ & $0.4737$ & 0 \\
  0.1 & 0.5 & $0.0271$ & $0.5920$ & $0.0161$ & $0.4695$ & 0 \\
  0.5 & 0.01 & $0.0251$ & $0.6000$ & $0.0151$ & $0.4969$ & 0 \\
  0.5 & 0.1 & $0.0244$ & $0.6093$ & $0.0149$ & $0.4919$ & 0 \\
  0.5 & 0.5 & $0.0254$ & $0.6107$ & $0.0155$ & $0.4835$ & 0 \\
  1.0 & 0.01 & $0.0232$ & $0.6053$ & $0.0140$ & $0.4947$ & 0 \\
  1.0 & 0.1 & $0.0274$ & $0.6027$ & $0.0165$ & $0.4842$ & 0 \\
  1.0 & 0.5 & $0.0258$ & $0.6147$ & $0.0159$ & $0.4842$ & 0 \\
%VICREG_DATA_END
\hline
\end{tabular}
\end{table}

Table~\ref{tab:vicreg} shows the VICReg grid search.
The best configuration ($\lambda_v=0.01, \lambda_c=0.1$) achieves NPMI $= 0.0279$ at seed=42,
a $+15\%$ gain over the baseline ($0.0242$), selected by NPMI$\times$TD (primary) with AU as tiebreaker.

Table~\ref{tab:vicreg_multiseed} validates this config across 3 seeds.
\textbf{Key finding}: VICReg \emph{reduces training variance} substantially
(20NG std: $0.0030 \to 0.0014$, $-53\%$; NYT: $0.0075 \to 0.0092$, mixed),
but does \emph{not} consistently improve mean NPMI
($0.0263$ vs.\ $0.0277$ on 20NG; $0.0336$ vs.\ $0.0372$ on NYT).
AU$=0$ persists across all seeds, confirming that the softmax bottleneck
cannot be resolved by $z$-space variance regularization alone.
The $+15\%$ at seed=42 reflects a favorable initialization rather than a robust improvement.
VICReg is therefore a \emph{training stabilizer}: useful for reducing seed sensitivity,
but not a reliable coherence booster.

\begin{table}[t]
\centering
\caption{Best VICReg config multi-seed results ($\lambda_v^*, \lambda_c^*$ from grid search),
3 seeds on 20NG and NYT ($K=50$). Compared to baseline EBM-LTM (no VICReg).}
\label{tab:vicreg_multiseed}
\small
\begin{tabular}{llrrrr}
\hline
\textbf{Dataset} & \textbf{Model} & \textbf{NPMI} & \textbf{TD} & \textbf{Acc} & \textbf{AU} \\
\hline
%VICREG_MULTI_DATA_START
  20NG & EBM-LTM (no VICReg) & $0.0277\pm0.0030$ & $0.616\pm0.003$ & $0.473\pm0.013$ & $0.0\pm0.0$ \\
  20NG & EBM-LTM + VICReg & $0.0263\pm0.0014$ & $0.6155\pm0.0111$ & $0.4737\pm0.0149$ & $0.0\pm0.0$ \\
  NYT & EBM-LTM (no VICReg) & $0.0372\pm0.0075$ & $0.646\pm0.011$ & $0.673\pm0.011$ & $0.0\pm0.0$ \\
  NYT & EBM-LTM + VICReg & $0.0336\pm0.0092$ & $0.6458\pm0.0079$ & $0.6765\pm0.0093$ & $0.0\pm0.0$ \\
%VICREG_MULTI_DATA_END
\hline
\end{tabular}
\end{table}

\subsection{Step Size Ablation}
\label{sec:stepsize_ablation}

\begin{table}[t]
\centering
\caption{ETM+Langevin step size ablation on 20NG ($K=50$, seed=42, $L_v=15$).
Exploration distance $\|z_{L_v} - z_0\|_2$ measures how far Langevin moves from the
amortized initializer. Optimal step size balances coherence and diversity.}
\label{tab:stepsize}
\small
\begin{tabular}{rrrrrr}
\hline
$\delta_v$ & \textbf{NPMI} & \textbf{TD} & \textbf{AU} & \textbf{Acc} & $\|z_{L_v}-z_0\|_2$ \\
\hline
%STEPSIZE_DATA_START
  0.01 & $0.0230$ & $0.5507$ & 0 & $0.3704$ & $2.838277$ \\
  0.02 & $0.0248$ & $0.6040$ & 0 & $0.4830$ & $4.076031$ \\
  0.05 & $0.0236$ & $0.6667$ & 0 & $0.5655$ & $6.468773$ \\
  0.1 & $0.0096$ & $0.7187$ & 0 & $0.5636$ & $9.043903$ \\
  0.2 & $0.0056$ & $0.7107$ & 0 & $0.5307$ & $12.731648$ \\
%STEPSIZE_DATA_END
\hline
\end{tabular}
\end{table}

Table~\ref{tab:stepsize} ablates the Langevin step size $\delta_v$.
The exploration distance $\|z_{L_v} - z_0\|_2$ increases monotonically with $\delta_v$
($2.84 \to 12.73$), but NPMI is not monotone:
$\delta_v{=}0.02$ achieves the best NPMI ($0.0248$), while larger steps
($\delta_v{=}0.1$: $0.0096$; $\delta_v{=}0.2$: $0.0056$) introduce MCMC noise
that corrupts the decoder signal, and smaller steps ($\delta_v{=}0.01$: $0.0230$)
fail to sufficiently explore beyond the encoder's mode.
The default $\delta_v = 0.02$ is the empirically optimal sweet spot.

\subsection{Held-Out Perplexity}
\label{sec:ppl}

\begin{table}[t]
\centering
\caption{Held-out perplexity ($K=50$, seed=42). PPL-amort: encoder mean $z_0$ (no Langevin).
PPL-refined: Langevin-refined $z_{L_v}$. Lower is better. BERTopic/FASTopic: no generative
model, hence not applicable (---).}
\label{tab:ppl}
\small
\begin{tabular}{llrr}
\hline
\textbf{Dataset} & \textbf{Model} & \textbf{PPL-amort} & \textbf{PPL-refined} \\
\hline
%PPL_DATA_START
\multicolumn{4}{l}{\textit{20NG}} \\
  & ETM & $1599.80$ & $---$ \\
  & ECRTM & $1416.41$ & $---$ \\
  & ETM+Langevin & $2270.95$ & $1906.58$ \\
  & EBMLTM & $2265.83$ & $1876.59$ \\
  & BERTopic & --- & --- \\
  & FASTopic & --- & --- \\
\hline
\multicolumn{4}{l}{\textit{NYT}} \\
  & ETM & $2545.13$ & $---$ \\
  & ECRTM & $2481.27$ & $---$ \\
  & ETM+Langevin & $3685.93$ & $2990.30$ \\
  & EBMLTM & $3684.57$ & $2983.46$ \\
  & BERTopic & --- & --- \\
  & FASTopic & --- & --- \\
\hline
\multicolumn{4}{l}{\textit{IMDB}} \\
  & ETM & $1459.38$ & $---$ \\
  & ECRTM & $1399.30$ & $---$ \\
  & ETM+Langevin & $1874.97$ & $1839.27$ \\
  & EBMLTM & $1873.40$ & $1837.69$ \\
  & BERTopic & --- & --- \\
  & FASTopic & --- & --- \\
\hline
%PPL_DATA_END
\end{tabular}
\end{table}

Table~\ref{tab:ppl} reports held-out perplexity.
PPL-refined is always lower than (or equal to) PPL-amortized for Langevin models,
confirming that posterior refinement improves the generative model's fit beyond
the amortized encoder.
ETM+Langevin and EBM-LTM are the only models with PPL-refined, providing
a unique evaluation axis that cluster-based models (BERTopic, FASTopic) cannot
provide --- reinforcing the generative modeling advantage of VAE-based approaches.

% ========================================================================
\section{Analysis}
\label{sec:analysis}
% ========================================================================

\subsection{Posterior Collapse and the AU=0 Problem}
\label{sec:collapse}

\textbf{All Langevin variants exhibit AU $= 0$} across all three datasets and all
AU regularizer weights tested. This indicates structural posterior collapse, not
a tuning issue.

\begin{table}[t]
\centering
\caption{Active-Unit regularizer ablation on 20NG ($K=50$). AU = 0 at all weights;
the regularizer modestly affects NPMI but cannot resolve the structural collapse.}
\label{tab:au_ablation}
\small
\begin{tabular}{llrrrr}
\hline
\textbf{Model} & \textbf{$\lambda_{\mathrm{AU}}$} &
    \textbf{NPMI} & \textbf{TD} & \textbf{Acc} & \textbf{AU} \\
\hline
ETM+Langevin & 0.0 & $0.0248$ & $0.604$ & $0.483$ & 0 \\
ETM+Langevin & 0.1 & $0.0276$ & $0.589$ & $0.480$ & 0 \\
ETM+Langevin & 0.5 & $0.0226$ & $0.605$ & $0.484$ & 0 \\
ETM+Langevin & 1.0 & $0.0260$ & $0.609$ & $0.494$ & 0 \\
\hline
EBM-LTM      & 0.0 & $0.0242$ & $0.616$ & $0.489$ & 0 \\
EBM-LTM      & 0.1 & $0.0222$ & $0.595$ & $0.481$ & 0 \\
EBM-LTM      & 0.5 & $0.0256$ & $0.603$ & $0.487$ & 0 \\
EBM-LTM      & 1.0 & $0.0276$ & $0.617$ & $0.500$ & 0 \\
\hline
\end{tabular}
\end{table}

Table~\ref{tab:au_ablation} shows results for the variance-penalizing regularizer
$\mathcal{L}_{\mathrm{AU}} = -\sum_k \min(\mathrm{Var}_x[\theta_k(x)], \tau)$
added to the main loss.
\textbf{AU remains 0 at all weights.} The regularizer modestly shifts NPMI
(best: ETM+Langevin at $\lambda=0.1$ achieves $0.0276$, a minor gain over $0.0248$),
but cannot resolve the collapse.

This is an important negative result. We identify three structural causes:
\begin{enumerate}
    \item \textbf{Softmax collapse}: $\theta = \softmax(z)$ maps large regions of
    $\R^K$ to near-uniform distributions. Small $z$-variance translates to
    near-zero $\theta$-variance.
    \item \textbf{Langevin mode convergence}: Small step size ($\delta_v = 0.02$)
    causes the chain to barely explore the posterior; samples cluster near the mode,
    producing low $\theta$-variance.
    \item \textbf{Imitation loss}: $\mathcal{L}_{\mathrm{init}} = \|z_0 - \sg(z_L)\|^2$
    drives the encoder toward a single-point estimate, reducing diversity.
\end{enumerate}

The variance regularizer cannot overcome cause (1): even with high $\lambda_{\mathrm{AU}}$,
the softmax bottleneck prevents $\theta$ from having variance $> 0.01$.

We address cause (2) via VICReg-style variance regularization on $z$-space
(Section~\ref{sec:vicreg}): by penalizing dimensions with below-target standard deviation
\emph{before} the softmax, the regularizer operates where the bottleneck is not yet active.
At seed=42, the best VICReg config achieves $+15\%$ NPMI (Table~\ref{tab:vicreg}).
However, multi-seed evaluation reveals mixed results: VICReg reduces training variance
by $53\%$ on 20NG (std: $0.0030 \to 0.0014$) but does not consistently improve
mean NPMI (Table~\ref{tab:vicreg_multiseed}).
AU$=0$ persists across all VICReg configurations: the softmax still compresses
$z$-variance into near-uniform $\theta$.
VICReg is therefore best understood as a \emph{stabilization} technique that reduces
seed-to-seed variance without reliably improving coherence.
Remaining avenues for future work: (i) replacing softmax with a Dirichlet-parameterized
reparameterization, (ii) Metropolis-Hastings correction for larger Langevin step sizes,
or (iii) explicit orthogonality constraints on topic embeddings.

\subsection{TD Decline and the Mode-Convergence Hypothesis}
\label{sec:td_decline}

Langevin refinement consistently reduces TD: from $0.853 \to 0.604$ on 20NG and
$0.827 \to 0.639$ on NYT. This is the same mechanism causing AU$=0$.

When Langevin SGLD is initialized at $z_0 \sim q_\lambda(z|x)$ with small step size
($\delta_v = 0.02$), the chain makes limited exploration and converges to the
local mode of $-\log p(z|x)$ closest to $z_0$.
If the posterior is approximately unimodal (which is likely for a BoW softmax model),
all documents converge to the same mode trajectory, producing:
(i) low variance across $\theta = \softmax(z)$ (AU$=0$), and
(ii) concentrated topic assignments (low TD).

The NPMI improvement despite AU$=0$ occurs because mode-convergence reduces
reconstruction variance---concentrating topics sharpens the top-$K$ words per topic,
which the NPMI metric rewards. This creates a fundamental trade-off:
\emph{Langevin improves coherence (NPMI) at the cost of diversity (TD)}.
The NPMI$\times$TD composite metric partially accounts for this
(ETM+Langevin: $0.015$ vs.\ ETM: $0.010$ on 20NG, a $+50\%$ gain).

The solution requires larger step sizes with Metropolis-Hastings correction or
annealing schedules that allow exploration early in training.

\paragraph{Unified explanation of three failure modes.}
We argue that AU$=0$, TD decline, and classification degradation are three manifestations
of a single underlying phenomenon: \emph{mode convergence} driven by small Langevin step size.

\begin{enumerate}
    \item \textbf{AU$=0$}: Langevin converges to a single posterior mode per document,
    producing low $z$-variance across the batch. The softmax bottleneck maps this
    concentrated $z$-distribution to near-uniform $\theta$, giving $\mathrm{Var}(\theta_k) < 0.01$.
    \item \textbf{TD decline}: Documents with similar content converge to similar modes,
    concentrating topic assignments and reducing the fraction of unique top-$K$ words.
    \item \textbf{Classification degradation}: When $\theta$ is near-uniform (AU$=0$),
    the SVM trained on $\theta$ loses its discriminative signal. The consistent ordering
    Acc$(\text{ETM}) >$ Acc$(\text{ETM+Langevin})$ across 20NG and NYT follows directly
    from the fact that AU$=0$ topics carry no class-discriminative information.
\end{enumerate}

This unified view explains why the AU regularizer fails: it targets symptom (1) but
cannot break the mode convergence that causes all three. VICReg (Section~\ref{sec:vicreg})
targets the root cause --- low batch variance in $z$-space --- and is empirically more
effective (Section~\ref{sec:vicreg_ablation}).

\subsection{Empirical Amortization Gap: Why Langevin Fails on IMDb}
\label{sec:amort_gap_empirical}

On IMDb, neither Langevin model improves over ETM: ETM achieves NPMI $= 0.0144$,
while ETM+Langevin scores $-0.0085$ and EBM-LTM scores $-0.0079$.
We now provide \emph{empirical} evidence, rather than speculation, for this failure.

\begin{table}[h]
\centering
\caption{Empirical amortization gap: how much Langevin refinement helps the encoder.
$\Delta$L2: mean Euclidean distance $\|z_L - z_0\|$ (step size).
$\Delta$logLik: mean log-likelihood improvement $\log p(x|z_L) - \log p(x|z_0)$ (benefit).
IMDb shows the same step size but $10\times$ smaller benefit, confirming
its encoder already approximates the posterior well.}
\label{tab:amort_gap}
\small
\begin{tabular}{lrr}
\hline
\textbf{Dataset} & \textbf{$\Delta$L2} & \textbf{$\Delta$logLik} \\
\hline
20NG  & $4.08 \pm 0.05$ & $\mathbf{18.65} \pm 10.7$ \\
NYT   & $4.33 \pm 0.06$ & $\mathbf{36.39} \pm 2.9$ \\
IMDb  & $3.89 \pm 0.03$ & $1.79 \pm 0.5$ \\
\hline
\end{tabular}
\end{table}

Table~\ref{tab:amort_gap} reveals that across all three datasets, Langevin moves
$z$ by a similar L2 distance ($\approx 4$). However, the log-likelihood benefit
of that move is $18.7\times$ on 20NG and $20.3\times$ on NYT compared to just $1.79$ on IMDb.

\textbf{Interpretation.} On 20NG and NYT, the encoder's initial $z_0$ is far from
the posterior mode in log-likelihood space---Langevin refinement corrects a
large inference error. On IMDb, $z_0$ is already near-optimal: the step moves $z$
by the same distance, but reconstruction barely improves ($\Delta$logLik $= 1.79$
vs.\ $18.65$ on 20NG). This confirms that IMDb's binary-class structure is simple
enough for the amortized encoder to approximate the posterior well, leaving no
room for Langevin to help.

This finding provides a principled criterion for when Langevin refinement is useful:
\emph{it helps when the amortization gap (measured by $\Delta$logLik) is large.}
Practitioners can estimate this gap cheaply on a small held-out set before
committing to full training with Langevin.

\subsection{Computational Overhead}

EBM-LTM costs 305s vs.\ ETM's 16s on 20NG ($19\times$). This is dominated by
the $L_v + L_s = 45$ sequential gradient evaluations per batch. Key reductions:
\begin{itemize}
    \item \textbf{ETM+Langevin} ($L_s = 0$): reduces to 198s ($12\times$ ETM)---roughly
    halving the EBM overhead while preserving most of the NPMI gain.
    \item \textbf{Reducing $L_v$}: $L_v = 10$ achieves NPMI $0.0228$ in 276s;
    $L_v = 5$ achieves $0.0202$ in 242s. Either is a viable efficiency-accuracy trade-off.
    \item \textbf{Future work}: normalizing flows \citep{rezende2015} could amortize the
    prior MCMC cost; adaptive step size schedules could reduce $L_v$ over training.
\end{itemize}

% ========================================================================
\section{Limitations}
\label{sec:limitations}
% ========================================================================

We enumerate our limitations explicitly:

\begin{enumerate}
    \item \textbf{Dataset scale}: Experiments are limited to three small-to-medium benchmarks
    (20NG: 11K docs; IMDb: 25K; NYT: 8K). Claims about generalization to larger corpora
    (Wiki, WOS) remain unverified.

    \item \textbf{Baseline coverage}: TopClus and other post-2022 hybrid models
    (combining clustering with LLMs) are not evaluated. BERTopic and FASTopic are
    included in our extended comparison (Table~\ref{tab:ext_baselines}), but the
    comparison is limited to offline settings without pretrained sentence encoders.

    \item \textbf{AU=0 collapse}: All Langevin variants exhibit persistent posterior
    collapse at the $\theta$-space level ($\mathrm{Var}(\theta_k) < 0.01$). The AU regularizer
    (Table~\ref{tab:au_ablation}) cannot resolve this due to the softmax bottleneck.
    Our VICReg-style regularization (Section~\ref{sec:vicreg}, Table~\ref{tab:vicreg})
    improves seed-42 NPMI by $+15\%$ and reduces training variance by $53\%$,
    but multi-seed mean NPMI is not consistently improved (Table~\ref{tab:vicreg_multiseed})
    and AU$=0$ persists, since the softmax still compresses $z$-variance into $\theta$.
    True resolution likely requires a Dirichlet reparameterization or Metropolis-Hastings correction.

    \item \textbf{Computational cost}: $12$--$19\times$ overhead over ETM limits
    applicability to large-scale settings.

    \item \textbf{Theoretical assumptions}: Theorem~\ref{thm:amort_gap} requires an LSI
    condition that may not hold for the non-log-concave topic model posterior
    (see Remark~\ref{rmk:gap}).

    \item \textbf{EBM prior's marginal contribution}: The EBM prior provides
    only marginal NPMI improvement over ETM+Langevin, raising the question of whether
    the additional complexity is justified in practice.
\end{enumerate}

% ========================================================================
\section{Conclusion}
\label{sec:conclusion}
% ========================================================================

We proposed EBM-LTM and systematically studied the effect of Langevin posterior
refinement and energy-based priors on neural topic modeling.
Our main finding is that \textbf{Langevin refinement is the dominant contributor}:
ETM+Langevin raises NPMI by $+116\%$ over ETM on 20NG ($0.0303\pm0.0049$ vs.\
$0.0139\pm0.0096$ across 3 seeds), confirming robustness beyond a single run.
The energy-based prior adds modest topic diversity ($+0.012$ TD) but does not
substantially improve coherence beyond ETM+Langevin across seeds
(NYT multi-seed: $0.0372\pm0.0075$ vs.\ $0.0371\pm0.0071$, indistinguishable within noise),
confirming that Langevin refinement is the essential component.

We provide five novel analytical contributions beyond the basic training recipe:
(1) empirical measurement of the amortization gap showing that Langevin's benefit is
dataset-dependent ($\Delta$logLik: $36.39$ on NYT, $18.65$ on 20NG, $1.79$ on IMDb),
cross-validated against Theorem~\ref{thm:amort_gap}'s prediction;
(2) a unified explanation of TD decline, AU$=0$, and classification degradation
as three manifestations of mode-convergence under small-step Langevin;
(3) VICReg-style variance regularization in $z$-space (Section~\ref{sec:vicreg})
that reduces training variance by $53\%$ (Table~\ref{tab:vicreg_multiseed}), though mean NPMI
is not consistently improved across seeds, with AU$=0$ persisting due to the softmax bottleneck;
(4) a step-size ablation connecting exploration distance $\|z_{L_v}-z_0\|_2$ to
topic quality tradeoffs (Table~\ref{tab:stepsize}); and
(5) held-out perplexity distinguishing PPL-amortized from PPL-refined, providing
an evaluation axis unavailable to cluster-based models (Table~\ref{tab:ppl}).

The factorial ablation provides a clean, reproducible demonstration of prior$\times$inference
effects. We hope the ETM+Langevin result---a simple, well-understood improvement---
motivates future work on Langevin refinement for topic models.

\paragraph{Future directions.}
(i) Apply Langevin with Metropolis-Hastings correction to resolve AU$=0$ more completely;
(ii) extend to larger corpora (Wiki, WOS) and LLM-based topic models;
(iii) learn adaptive Langevin schedules to reduce $12$--$19\times$ overhead;
(iv) investigate whether the amortization gap criterion ($\Delta$logLik $< 5$ threshold)
generalizes to other VAE-based generative models.

% ========================================================================
\section*{Acknowledgements}
% ========================================================================

We thank the TopMost team for their excellent open-source framework.

% ========================================================================
\bibliographystyle{abbrvnat}
\bibliography{ebmltm_refs}
% ========================================================================

% ========================================================================
\appendix
% ========================================================================

\section{Full Proof of Theorem~\ref{thm:amort_gap}}
\label{app:proof}

\paragraph{Setup.}
Let $\pi = p(\cdot|x)$ and $U(z) = -\log\pi(z)$.
Assume $\pi$ satisfies the log-Sobolev inequality (LSI) with constant $\rho > 0$:
for any density $q$ with $q \ll \pi$,
\[
    \mathrm{KL}(q\|\pi) \;\le\; \frac{1}{2\rho}\,\mathcal{I}(q\|\pi),
    \quad \mathcal{I}(q\|\pi) = \E_q\!\left[\left\|\nabla\log\frac{q}{\pi}\right\|^2\right].
\]
Assume $\nabla \log\pi$ is $M$-Lipschitz.

\paragraph{Step 1: KL decay under the continuous Langevin SDE.}
The overdamped Langevin SDE $dz = -\nabla U(z)\,dt + \sqrt{2}\,dW$ has stationary
distribution $\pi$. Under the LSI, the relative entropy decays exponentially:
$\frac{d}{dt}\mathrm{KL}(q_t\|\pi) = -\mathcal{I}(q_t\|\pi) \le -2\rho\,\mathrm{KL}(q_t\|\pi)$.
Gronwall's inequality gives $\mathrm{KL}(q_t\|\pi) \le e^{-2\rho t}\,\mathrm{KL}(q_0\|\pi)$.

\paragraph{Step 2: KL to $W_2$ via Talagrand's inequality.}
The LSI implies Talagrand's $T_2$ inequality \citep{otto2000generalization}:
$W_2(q,\pi)^2 \le \frac{2}{\rho}\,\mathrm{KL}(q\|\pi)$. Combined with Step 1:
$W_2(q_t,\pi) \le e^{-\rho t}\,W_2(q_0,\pi) \le e^{-\rho t}\,\varepsilon$.

\paragraph{Step 3: Euler--Maruyama discretization.}
Each discrete step accumulates a bias bounded by $C_1 M\sqrt{\delta K}$ in $W_2$
\citep{dalalyan2017theoretical}. After $T$ steps with accumulated continuous
decay $e^{-\rho T\delta}$ and $T$ bias increments (summed geometrically):
\[
    W_2(q_T,\pi) \;\le\; e^{-\rho T\delta/2}\,\varepsilon + \frac{C M}{\rho}\sqrt{\delta K}.
\]
\qed

\section{Implementation Details}
\label{app:impl}

\paragraph{Energy network.}
SN-Linear($K$, $2K$) $\to$ SiLU $\to$ SN-Linear($2K$, $2K$) $\to$ SiLU $\to$ SN-Linear($2K$, 1),
where SN denotes spectral normalization via power iteration.

\paragraph{Encoder.}
Linear($V$, 512) $\to$ ReLU $\to$ Linear(512, 512) $\to$ ReLU $\to$ Dropout(0)
$\to$ \{Linear(512, $K$) for $\mu$, Linear(512, $K$) for $\log\sigma^2$\}.
Input: $\ell_1$-normalized BoW.

\paragraph{Decoder (ETM-style).}
$\beta_{kv} = \text{softmax}_v(\rho_v^\top \alpha_k)$; pretrained word embeddings
$\rho$ frozen; topic embeddings $\alpha$ trained.
Reconstruction: $\hat x = \theta^\top \beta$, loss $= -\sum_v x_v \log \hat x_v$.

\paragraph{MCMC efficiency.}
Energy parameters frozen (\texttt{requires\_grad\_(False)}) during Langevin loops.
$\beta$ precomputed once per batch (not inside the loop).

\paragraph{Hyperparameter table.}
\begin{center}
\small
\begin{tabular}{lll}
\hline
Parameter & Value & Rationale \\
\hline
$K$ (topics) & 50 & Standard setting \\
embed dim $L$ & 200 & Match ETM \citep{dieng2020} \\
$L_v$ & 15 & Ablation optimum \\
$\delta_v$ & 0.02 & Stable with SN \\
$L_s$ & 30 & Short-run MCMC \citep{nijkamp2019shortrun} \\
$\delta_s$ & 0.1 & Gaussian base \\
Warmup & 25\% epochs & Prevent early collapse \\
$\eta$ & $2\times10^{-3}$ & Adam (encoder/decoder) \\
$\eta_\alpha$ & $10^{-4}$ & Adam (energy, stable CD) \\
Batch & 200 & Standard \\
Grad clip & 5.0 (main), 1.0 (energy) & Prevent divergence \\
\hline
\end{tabular}
\end{center}

\paragraph{Missing baselines.}
BERTopic \citep{grootendorst2022} requires a sentence-transformer backend not
integrated with TopMost; FASTopic \citep{wu2024fastopic} and TopClus
\citep{meng2022topclus} were not available in our experimental setup.
We acknowledge this is a significant limitation and list it in Section~\ref{sec:limitations}.

\end{document}
